\documentclass[10pt]{article}
\usepackage[margin=20mm]{geometry}
\usepackage{amsmath,amssymb,mathtools,fontspec,unicode-math,xurl,hyperref}
\IfFontExistsTF{DejaVu Serif}{\setmainfont{DejaVu Serif}}{\setmainfont{Latin Modern Roman}}
\setmathfont{latinmodern-math.otf}
\setlength{\parindent}{0pt}\setlength{\parskip}{6pt}\sloppy
\title{Initial scientific review of
Étale cohomology of algebraic varieties over Stein compacta}
\author{}
\date{26 October 2026}
\begin{document}
\maketitle
Abstract: A concrete normal-crossing Stein surface shows that the fiber product prescribed in the normalization reduction of Proposition 4.2 need not kill the cohomology class. An extra component, arising from two different normalization branches, retains a nonzero class in degree one with coefficients in \(\mathbb Z/2\). Taking the product over the normalization repairs this local construction. This is a Major defect in a proof step used by the main comparison theorem, not a counterexample to that theorem. A separate Minor scope defect concerns the ambient-dimension equality stated in Theorem 0.4; the body already states the appropriate pure-dimension hypothesis. The review concerns only the acquired arXiv v3 manuscript. Independent scientific acceptance remains pending.\par
\section*{Identity, version and inspected scope}
The assigned paper is Olivier Benoist, Étale cohomology of algebraic varieties over Stein compacta, M25-0063. The exact reviewed file is arXiv:2305.06054v3, dated 3 February 2025, with 36 PDF pages. Printed page numbers and one-based PDF page numbers agree. The acquired PDF SHA-256 is\par
The associated journal record is Inventiones mathematicae 240 (2025), 497–536, DOI 10.1007/s00222-025-01323-2 [paper, publisher]. The publisher version of record was not acquired. The arXiv record’s description of its version as final does not establish file or content identity with the journal article. None is asserted here.\par
All 1,939 lines of the native text extraction were read, including references; truncated display segments were retrieved again. Nineteen original PDF pages were also inspected visually. The scientific review concentrated on the hypotheses and dependency chain of Propositions 2.4, 4.1–4.3, Theorems 3.9, 5.1, 5.5–5.10, Corollary 5.11, Remark 5.12 and the sums-of-squares applications. Reading the text is not an independent reproof of every cited theorem. The released TeX was inspected statically at the disputed formulas and the corresponding body statements; it was not compiled or executed. No separate essential supplement was identified. Optional software absence is not a defect.\par
\section*{Definitions, quantifiers and proof dependencies}
A Stein compact subset has a basis of Stein open neighborhoods. \(\mathcal O(K)\) and \(\mathcal M(K)\) are rings of germs on neighborhoods of \(K\); they are not rings of functions merely on the underlying point set. In §1.1 the paper permits general second-countable complex spaces, without an implicit connectedness or pure-dimension assumption. Normality is local and does not force distinct connected components to have the same dimension.\par
Theorem 5.1, pp. 26–28, compares cohomology only after taking the directed limit over Stein neighborhoods of \(K\). The algebraic scheme is interiorly of finite type and the sheaf is torsion; either interior properness or interior constructibility is required. Corollary 5.2 removes the neighborhood-limit presentation under properness. These alternatives cannot both be dropped: the source’s Remark 3.10 explicitly gives a counterexample.\par
The checked dependency chain is as follows. Proposition 2.4 kills a positive-degree class after a finite, possibly ramified, cover, using a degree-one cyclic cover and a Mayer–Vietoris induction through bumps. Proposition 4.1 uses the relative comparison Theorem 3.9 and a finite-cover annihilation theorem to obtain local vanishing on a normal space. Proposition 4.2 extends this to general Stein spaces and all points over a chosen base point. Proposition 4.3 globalizes that result so a ramified class extends. Theorem 5.1, Step 5, combines Propositions 4.3 and 2.4 to kill higher direct images in the qfh comparison. Problem 1 concerns the first reduction in Proposition 4.2, before its separate Galois-closure argument.\par
For the equivariant application, \(G=\operatorname{Gal}(\mathbb C/\mathbb R)\) acts antilinearly. Theorem 5.5 uses a comparison of Hochschild–Serre spectral sequences. The absence of fixed points in Theorem 5.6 is essential: it permits the equivariant Artin vanishing input. The statement of that external input was checked in [Retained latex: \textbackslash{}cite[Theorem\textasciitilde{}2.6]\{vanishing\}]; taking its compact subset empty and using an empty fixed locus yields the vanishing used here. Its entire proof was not independently reproduced.\par
\section*{Problem 1: the product in the normalization reduction has extra components}
\subsection*{Source claim and location}
Proposition 4.2 is on pp. 23–24. Its proof starts with the normalization \(\nu:\widetilde S\to S\), the points \(\widetilde s_i\) above \(s\), and finite surjections \(p_i:T_i\to\widetilde S\) which kill \(\nu^*\alpha\) over disjoint neighborhoods \(\Omega_i\) of \(\widetilde s_i\). The first sentence on p. 24 asserts that the fiber product over \(S\) of the maps \(\nu\circ p_i\) solves the original problem. The same construction is explicit in the released `compactStein.tex`, line 1299 (one-based). The essential inference is \[p_i^*\nu^*\alpha=0\ \hbox{over }p_i^{-1}(\Omega_i\setminus\nu^{-1}Z)
 \quad\Longrightarrow\quad
 p^*\alpha=0\ \hbox{near }p^{-1}(s)\setminus p^{-1}Z\] for \(p:\prod_S T_i\to S\). It need not hold.\par
\subsection*{Explicit parameters and verification of every local assumption}
Let \[\begin{aligned}
 S&=\{(x,y,z)\in\mathbb C^3:xy=0\},&
 A&=\{y=0\},\quad B=\{x=0\},\\
 s&=(0,0,0),&
 U&=S,\quad Z=\{z=0\}\cap S .
\end{aligned}\] This is a reduced Stein space, being a closed analytic subspace of \(\mathbb C^3\); \(U\) is open in \(S\), \(s\in U\), and \(Z\) is a closed analytic subset. The two planes \(A\) and \(B\) meet in the \(z\)-axis. The normalization is the finite surjection \[\nu:\widetilde S=A\amalg B\longrightarrow S.\] It has precisely two points over \(s\). Take \(k=1\) and \(m=2\). Let \(\eta\in H^1(\mathbb C^*,\mathbb Z/2)\) be the nonzero class, and define \[\alpha=z^*\eta\in H^1(S\setminus Z,\mathbb Z/2).\] The inclusion of the punctured \(z\)-axis is a section of \(z:S\setminus Z\to\mathbb C^*\), so \(\alpha\) is nonzero.\par
Define two covers of \(\widetilde S\): \[\begin{array}{c|cc}
 &\text{over }A&\text{over }B\\ \hline
p_A:T_A\to\widetilde S &(u,t)\mapsto(u,0,t^2)&(v,z)\mapsto(0,v,z)\\
p_B:T_B\to\widetilde S &(u,z)\mapsto(u,0,z)&(v,t)\mapsto(0,v,t^2).
\end{array}\] Each source is a disjoint union of two copies of \(\mathbb C^2\). Each map is holomorphic, finite and surjective, since its component maps are a degree-two polynomial map and an identity. Thus these maps are allowed inputs to the reduction.\par
Choose \(\Omega_A=A\) and \(\Omega_B=B\), viewed as disjoint open subsets of the normalization. They are neighborhoods of the respective preimages of \(s\); their union is all of \(\widetilde S\). On \(p_A^{-1}(\Omega_A\setminus\nu^{-1}Z)\), the map on the punctured coordinate is \(t\mapsto t^2\), so \[p_A^*\nu^*\alpha=2\eta=0\quad\text{in }H^1(-,\mathbb Z/2).\] The same calculation gives \(p_B^*\nu^*\alpha=0\) over \(\Omega_B\). Hence every input vanishing condition asserted in the reduction holds. For every neighborhood \(\Omega\) of \(s\) one also has \(\nu^{-1}(\Omega)\subset\Omega_A\cup\Omega_B\).\par
\subsection*{The prescribed output and its actual cohomology}
Form the prescribed product \[T=T_A\times_S T_B,\qquad p:T\to S .\] It is again finite and surjective. However, take the open-and-closed part in which the first factor is the identity copy over \(B\) and the second is the identity copy over \(A\). This part is \[D=B\times_S A=\{x=y=0\}\simeq\mathbb C_z.\] The map \(p|_D\) is the identity onto the intersection axis. Consequently \[(p^*\alpha)|_{D\setminus p^{-1}Z}=\eta\ne0.\] This obstruction persists after every shrinking of \(\Omega\) around \(s\): a sufficiently small circle on the punctured \(z\)-axis lies in \(\Omega\) and has evaluation \(1\) against the restricted mod-two class. Thus the claimed vanishing is false for this permitted product. It is not a failure to choose a sufficiently small neighborhood.\par
An independent finite calculation of cyclic covering degrees confirms the four branch-pair cases. For branch pair \(AA,AB,BA,BB\), respectively, the degrees of the two coordinate covers are \((2,1),(2,2),(1,1),(1,2)\). Over a punctured base, their fiber product has \(\gcd(d,e)\) components, each with winding \(\operatorname{lcm}(d,e)\). The respective winding parities are \[(0,0,1,0).\] Only the \(BA\) component already proves the contradiction. The analytic construction above, rather than this finite calculation, verifies the Stein and neighborhood hypotheses.\par
\subsection*{Counterevidence, interpretation and core impact}
The proposition’s existence statement is not disproved: there are covers which kill the class, including the corrected product below. Nor is the normal-space Galois-closure argument in the second paragraph of the proof contradicted. The issue is the explicit reduction from a general Stein space to its normalization.\par
Taking products over \(\widetilde S\) would force the factors to have the same normalization lift. That is a different mathematical object from the stated product of \(\nu\circ p_i\) over \(S\). The example verifies a nonzero cohomology class on an actual additional component, rather than merely detecting an ambiguous symbol or an OCR error. The original PDF and the matching released TeX both specify the latter construction.\par
The supported severity is Major. Proposition 4.2 is the local input explicitly invoked in Proposition 4.3, p. 24, and Proposition 4.3 is explicitly invoked in Theorem 5.1, Step 5, p. 28. The written general-space proof of the main comparison theorem therefore needs this reduction corrected. This conclusion does not refute the comparison isomorphism, the dimension bounds, or the sums-of-squares theorem. An easy local repair does not remove the demonstrated dependency.\par
\subsection*{Precise local repair and its boundary}
Replace the prescribed product by \[T'=\prod_{\widetilde S}T_i,\qquad
 \widetilde p:T'\to\widetilde S,\qquad p'=\nu\circ\widetilde p.\] This is finite and surjective. Because \(\nu\) is finite and the disjoint neighborhoods \(\Omega_i\) contain the whole fiber over \(s\), choose \(\Omega\) around \(s\) with \(\nu^{-1}(\Omega)\subset\bigcup_i\Omega_i\). The subsets \(\widetilde p^{-1}(\Omega_i\cap\nu^{-1}\Omega)\) are disjoint and cover \((p')^{-1}\Omega\). On the \(i\)th such subset, the projection \(T'\to T_i\) pulls back the already vanishing class. Cohomology of a finite disjoint union is the product of the component cohomology groups. Therefore \((p')^*\alpha\) vanishes on \((p')^{-1}(\Omega\setminus Z)\), for every degree \(k\) in the proposition.\par
This proves the corrected normalization reduction in its stated generality. In the example it retains only the \(AA\) and \(BB\) parts, whose winding parities are both zero. It does not independently certify all subsequent Cech refinement arguments, hypercover descent, or external analytic inputs. Author adoption has not been requested or observed.\par
\section*{Problem 2: ambient dimension in the introductory equality}
\subsection*{Exact statement and counterexample}
Theorem 0.4 on p. 3 asserts \(\operatorname{cd}\mathcal M(K)=n\) for a connected Stein compact subset \(K\) of a normal Stein space \(S\) of dimension \(n\). The abstract on p. 1 has the same dimensional scope. In contrast, Remark 5.12(ii), p. 31, explicitly assumes that \(S\) has pure dimension \(n\) and \(K\) is nonempty.\par
Under the conventions of §1.1, take \[S=\{\ast\}\amalg\mathbb C,\qquad K=\{\ast\},\qquad n=1.\] This is a normal, second-countable Stein space of dimension one. The compact set \(K\) is nonempty and connected, and is Stein because it is itself an open Stein component and hence has a basis of Stein neighborhoods. These verify every stated assumption of Theorem 0.4. But restriction to that open component gives \[\mathcal O(K)=\mathcal M(K)=\mathbb C,\qquad \operatorname{cd}\mathcal M(K)=0,\] since \(\mathbb C\) is algebraically closed and has trivial absolute Galois group. The displayed conclusion would instead give \(1\). Adding an unrelated component changes \(\dim S\) without changing the germ along \(K\).\par
\subsection*{Counterevidence, impact and correction}
Normality does not exclude this example, and the theorem states connectedness of \(K\), not of \(S\). The important counterevidence is supplied by the paper itself: Corollary 5.11 only asserts an upper bound, and Remark 5.12(ii) already states the pure-dimension assumption needed for its lower bound. The above example satisfies the upper bound \(0\le1\) and does not meet the remark’s pure-dimension hypothesis.\par
This is a Minor scope defect in the introductory restatement. There is no demonstrated defect in the body result under its explicit hypotheses. The sums-of-squares bound in Theorem 6.4 uses the upper bound, so this example does not damage that quantitative application. The correction is to include nonemptiness of \(K\) and pure dimension \(n\) in Theorem 0.4, following Remark 5.12(ii), or to state the dimension of the normal connected component containing \(K\). The first formulation is already supported by the author’s body argument; it is not a new result claimed by this report.\par
\section*{Other checked reasoning and limits}
The Mayer–Vietoris step in Proposition 2.4 was checked at the level of the exact sequence: the two restrictions of the degree-\(k\) class vanish, so it lifts to a degree-\((k-1)\) class on the overlap; the induction hypothesis kills that lift on a finite cover, and naturality kills the original class after shrinking. Holomorphic convexity of the pulled-back compact set is supplied by Proposition 1.6(i). The bump construction and the cited approximation theorems were read but not independently reproved.\par
The relative comparison Theorem 3.9 preserves the properness/constructibility alternatives and uses dimension induction after reduction to regular closed strata. The terminal local calculation is the cohomology of a punctured complex normal slice, with nonzero degrees \(0\) and \(2c-1\), matched to purity. No contrary calculation was established in that inspected reduction.\par
For Theorem 5.1, Step 5, p. 28, the factorization is printed with finite target map to \(Y_U\), while the following inverse image of \(K\subset U\) and Steinness argument require a finite map to \(X_U=\operatorname{Spec}\mathcal O(U)\). A precise consistent reading is to factor the composite \(W\to X_U\) by Zariski’s Main Theorem, which is legitimate because it is separated and quasi-finite of finite presentation. Then \(\overline W^{\mathrm{an}}\) is finite over the Stein space \(U\), the inverse image of \(K\) is compact and holomorphically convex, and the finite cover restricted to \(W\) still refines the original cover of \(Y_U\). This restores the stated uses of Propositions 4.3 and 2.4 at this point. This forced base-label clarification is recorded without a separate severity count; the report does not claim the literal label already provides the required compactification. Unlike Problem 1, no additional counterexample to a separately specified gluing output is being counted here.\par
Theorem 5.6 requires an affine scheme and empty real fixed locus. Passing to the function field in Theorem 5.7 uses distinguished affine opens and a filtered colimit, not an assertion that arbitrary nonaffine schemes have the same bound. In Proposition 6.3 the descending argument is also consistent at the endpoint: to descend to the symbol of degree \(n+1\), injectivity uses vanishing in degree \(n+1\) over \(F(\sqrt{-1})\), which follows from the bound \(\le n\). The split quadratic algebra and characteristic-two cases do not furnish a counterexample. Theorems 6.4–6.5 retain compactness and their local-ring hypotheses; no extension to every noncompact real-analytic space is inferred.\par
The limited exact-title, publisher and arXiv checks did not identify a specific erratum addressing the two findings. This is not an exhaustive absence claim. No newer or journal text was substituted for the pinned manuscript.\par
\section*{Disposition}
Item | Initial-review conclusion\par
Problem 1 | Major: the stated normalization product need not kill the class.\par
Repair 1 | Product over the normalization; local vanishing proof verified analytically.\par
Problem 2 | Minor: introductory equality lacks the body’s pure-dimension restriction.\par
Repair 2 | Use the hypotheses already stated in Remark 5.12(ii).\par
Core conclusions | Not refuted; the written general-space derivation needs correction.\par
Remaining work | Independent acceptance; full downstream proof and external input verification.\par
Adoption | Not requested or observed; body hypothesis for Problem 2 pre-exists this review.\par
Scientific decision | Pending; no human review recorded.\par
\section*{Supporting checks}
The source was extracted by the existing native Poppler command `pdftotext -layout`. Its extraction hash is\par
Original PDF pages actually viewed are 1, 3, 5, 14, 15, 19–24 and 26–33. The original diagrams on pp. 21 and 23 and the formulas essential to both findings were inspected as images. The ledger records exact text and static-TeX ranges. The complete manifest and all supplied source-file hashes were checked separately; indexing and hash verification are not scientific review.\par
Its two files are\par
> `normalization\_cover\_check.py`
> `normalization\_cover\_check.json`.\par
The existing native Python executable was invoked with that script as its sole program argument. It returned four branch-pair rows with winding parities \((0,0,1,0)\), a surviving component over \(S\), and zero parities for the two same-branch components retained over \(\widetilde S\). No author program was run. The field computation in Problem 2 is analytic and needs no numerical experiment.\par
The search command `rg` was unavailable; a read-only native text search was used instead. Some large tool displays were truncated and their relevant segments were reread before claiming coverage. These operational events do not supply scientific evidence.\par
\section*{Supporting checks}
This report is compiled from its exact saved TeX by the explicitly authorized portable Tectonic 0.17.0 in untrusted mode. There is no shell escape, installation or author-code execution. The cache snapshot is a superset of used bundled dependencies; the actual TeX log records inputs. Every final report page must be inspected before delivery. Nonfatal Fontconfig output is preserved in the build log.\par
The account context is the previously confirmed Minhao Liu’s Group—Academic Researchers. This assistant performed scientific reasoning, source inspection, calculation, image review and compilation; Work only schedules. The fields `human\_reviews=[]` and `scientific\_decision=pending` remain unchanged. Complete initial delivery is not independent scientific acceptance.\par
9 O. Benoist, Étale cohomology of algebraic varieties over Stein compacta, arXiv:2305.06054v3, 3 February 2025. Exact reviewed PDF and matching released TeX; hash above. Publisher metadata for the same title, Invent. Math. 240 (2025), 497–536, DOI 10.1007/s00222-025-01323-2; arXiv record https://arxiv.org/abs/2305.06054, checked 26 October 2026. Metadata only; journal PDF equivalence unestablished. O. Benoist, On the Artin vanishing theorem for Stein spaces, arXiv:2410.04291v1, 5 October 2024, Theorem 2.6. Consulted only for the statement and application range of the cited dependency.\par
\end{document}
